%
\chapter{FAQ}
\label{appdxa-faq} 

\section{FAQ: Answers to Frequently Asked Questions} 

Answers to frequently asked questions (FAQ) are compiled, summarized and
displayed in the following table. \textbf{\emph{Note:}} As more FAQ and
answers become available, they will be appended in the table for future
release.

\begin{longtable}[]{@{}|l|l|l|l|@{}}
\toprule
\begin{minipage}[b]{0.08\columnwidth}\raggedright
No.\strut
\end{minipage} & \begin{minipage}[b]{0.30\columnwidth}\raggedright
Question\strut
\end{minipage} & \begin{minipage}[b]{0.33\columnwidth}\raggedright
Answer\strut
\end{minipage} & \begin{minipage}[b]{0.17\columnwidth}\raggedright
Reference\strut
\end{minipage}\tabularnewline
\midrule
\endhead
 \begin{minipage}[t]{0.08\columnwidth}\raggedright
1\strut
\end{minipage} & \begin{minipage}[t]{0.30\columnwidth}\raggedright
I want to change the buzzer operation in the event of the occurrence of
an ``alarm or warning.''\strut
\end{minipage} & \begin{minipage}[t]{0.33\columnwidth}\raggedright
ON/OFF switch can be controlled via {[}Setup{]} → {[}Configuration{]} →
{[}Set Sound{]}. If the beep is turned off, alarm and warning are
indicated only with red notifications on the HMI.\strut
\end{minipage} & \begin{minipage}[t]{0.17\columnwidth}\raggedright
2.8 Setting the Alarm and Warning Buzzer\strut
\end{minipage}\tabularnewline
 \begin{minipage}[t]{0.08\columnwidth}\raggedright
2\strut
\end{minipage} & \begin{minipage}[t]{0.30\columnwidth}\raggedright
What should I do when a system abnormality has occurred?\strut
\end{minipage} & \begin{minipage}[t]{0.33\columnwidth}\raggedright
Notify the sales representative or customer support for
assistance.\strut
\end{minipage} & \begin{minipage}[t]{0.17\columnwidth}\raggedright
\strut
\end{minipage}\tabularnewline
 \begin{minipage}[t]{0.08\columnwidth}\raggedright
3\strut
\end{minipage} & \begin{minipage}[t]{0.30\columnwidth}\raggedright
I want to turn off the HMI immediately.\strut
\end{minipage} & \begin{minipage}[t]{0.33\columnwidth}\raggedright
To turn off the HMI immediately, press {[}Setup{]} → {[}Configuration{]}
→ {[}Display Setup{]}, then press \textbf{Turn Off Display}. To turn it
back on, tab your finger on the screen.\strut
\end{minipage} & \begin{minipage}[t]{0.17\columnwidth}\raggedright
7.9.1 Display Setup on HMI\strut
\end{minipage}\tabularnewline
 \begin{minipage}[t]{0.08\columnwidth}\raggedright
4\strut
\end{minipage} & \begin{minipage}[t]{0.30\columnwidth}\raggedright
I want to set the HMI in sleep mode.\strut
\end{minipage} & \begin{minipage}[t]{0.33\columnwidth}\raggedright
Several options are available for setting the HMI in sleep mode. Press
{[}Setup{]} → {[}Configuration{]} → {[}Display Setup{]}, then select
\textbf{HMI Screensaver timeout}. To turn it back on, tab your finger on
the screen.\strut
\end{minipage} & \begin{minipage}[t]{0.17\columnwidth}\raggedright
7.9.1 Display Setup on HMI\strut
\end{minipage}\tabularnewline
 \begin{minipage}[t]{0.08\columnwidth}\raggedright
5\strut
\end{minipage} & \begin{minipage}[t]{0.30\columnwidth}\raggedright
``OFF'' is displayed for a setup value.\strut
\end{minipage} & \begin{minipage}[t]{0.33\columnwidth}\raggedright
When the equipment is stopped or the humidity control is set to OFF,
``OFF'' is displayed.\strut
\end{minipage} & \begin{minipage}[t]{0.17\columnwidth}\raggedright
2.3.3 Monitor\strut
\end{minipage}\tabularnewline
 \begin{minipage}[t]{0.08\columnwidth}\raggedright
6\strut
\end{minipage} & \begin{minipage}[t]{0.30\columnwidth}\raggedright
``---'' is displayed as the temperature (humidity) process value.\strut
\end{minipage} & \begin{minipage}[t]{0.33\columnwidth}\raggedright
If the temperature (humidity) process value is abnormal or the humidity
control is set to OFF, ``---'' is displayed.\strut
\end{minipage} & \begin{minipage}[t]{0.17\columnwidth}\raggedright
2.3.3 Monitor\strut
\end{minipage}\tabularnewline
 \begin{minipage}[t]{0.08\columnwidth}\raggedright
7\strut
\end{minipage} & \begin{minipage}[t]{0.30\columnwidth}\raggedright
I want to change the refrigeration setup from Auto to Manual (arbitrary
state).\strut
\end{minipage} & \begin{minipage}[t]{0.33\columnwidth}\raggedright
Access {[}Set Mode{]}, select a Constant setup, then apply refrigeration
manual setting as needed. When you want to change the mode from Auto to
Manual during operation, make the manual setting first.\strut
\end{minipage} & \begin{minipage}[t]{0.17\columnwidth}\raggedright
4.3 Refrigeration Setting\strut
\end{minipage}\tabularnewline
 \begin{minipage}[t]{0.08\columnwidth}\raggedright
8\strut
\end{minipage} & \begin{minipage}[t]{0.30\columnwidth}\raggedright
When is the setting of time enabled?\strut
\end{minipage} & \begin{minipage}[t]{0.33\columnwidth}\raggedright
At the moment when the time signal is set, the contact output in
question changes Contact rating: 250VAC, 5A\strut
\end{minipage} & \begin{minipage}[t]{0.17\columnwidth}\raggedright
\strut
\end{minipage}\tabularnewline
 \begin{minipage}[t]{0.08\columnwidth}\raggedright
9\strut
\end{minipage} & \begin{minipage}[t]{0.30\columnwidth}\raggedright
Is the trend graph data erased when the circuit breaker is turned off or
when power failure occurs?\strut
\end{minipage} & \begin{minipage}[t]{0.33\columnwidth}\raggedright
Data are stored in the non-volatile memory of the GL controller system;
they cannot not be erased by power failure or shutdown.\strut
\end{minipage} & \begin{minipage}[t]{0.17\columnwidth}\raggedright
4.9.7 Monitor: Trend; 4.9.8 Trend Graph Manipulation Buttons\strut
\end{minipage}\tabularnewline
 \begin{minipage}[t]{0.08\columnwidth}\raggedright
10\strut
\end{minipage} & \begin{minipage}[t]{0.30\columnwidth}\raggedright
What is the ramp operation?\strut
\end{minipage} & \begin{minipage}[t]{0.33\columnwidth}\raggedright
The ramp operation is to control the temperatures from the set point of
the previous step to that of the present step at a certain ramp.\strut
\end{minipage} & \begin{minipage}[t]{0.17\columnwidth}\raggedright
5.3.1 Step-Wise Programming\strut
\end{minipage}\tabularnewline
 \begin{minipage}[t]{0.08\columnwidth}\raggedright
11\strut
\end{minipage} & \begin{minipage}[t]{0.30\columnwidth}\raggedright
What is the soak time control?\strut
\end{minipage} & \begin{minipage}[t]{0.33\columnwidth}\raggedright
After the operation of the step is started, the equipment waits until
the process value attains the attainment range of the set point, and
then starts the counting of time. Thus, the soak time in the temperature
set point is the same as the setup time.\strut
\end{minipage} & \begin{minipage}[t]{0.17\columnwidth}\raggedright
5.3.1 Step-Wise Programming\strut
\end{minipage}\tabularnewline
 \begin{minipage}[t]{0.08\columnwidth}\raggedright
12\strut
\end{minipage} & \begin{minipage}[t]{0.30\columnwidth}\raggedright
Can the soak time and ramp control be set simultaneously?\strut
\end{minipage} & \begin{minipage}[t]{0.33\columnwidth}\raggedright
The ramp control and soak time cannot be set simultaneously.\strut
\end{minipage} & \begin{minipage}[t]{0.17\columnwidth}\raggedright
5.3.1 Step-Wise Programming\strut
\end{minipage}\tabularnewline
 \begin{minipage}[t]{0.08\columnwidth}\raggedright
13\strut
\end{minipage} & \begin{minipage}[t]{0.30\columnwidth}\raggedright
Can the program in operation be edited or saved?\strut
\end{minipage} & \begin{minipage}[t]{0.33\columnwidth}\raggedright
During program execution, that program can be loaded into the program
editor for editing, but the edited version cannot be saved back in the
same slot that was loaded from for execution; it can however be saved in
a different slot (and under a different name).\strut
\end{minipage} & \begin{minipage}[t]{0.17\columnwidth}\raggedright
5.6.4 Manipulating a Program during its Execution\strut
\end{minipage}\tabularnewline
 \begin{minipage}[t]{0.08\columnwidth}\raggedright
14\strut
\end{minipage} & \begin{minipage}[t]{0.30\columnwidth}\raggedright
Can the program in operation be deleted?\strut
\end{minipage} & \begin{minipage}[t]{0.33\columnwidth}\raggedright
During program execution, that program cannot be deleted.\strut
\end{minipage} & \begin{minipage}[t]{0.17\columnwidth}\raggedright
5.6.4 Manipulating a Program during its Execution\strut
\end{minipage}\tabularnewline
 \begin{minipage}[t]{0.08\columnwidth}\raggedright
15\strut
\end{minipage} & \begin{minipage}[t]{0.30\columnwidth}\raggedright
Can the program contain loops (counters) to reduce the number of
steps?\strut
\end{minipage} & \begin{minipage}[t]{0.33\columnwidth}\raggedright
Program may contain multiple loops (i.e., counters) to repeat steps or
groups of steps.\strut
\end{minipage} & \begin{minipage}[t]{0.17\columnwidth}\raggedright
5.3.1 Step-Wise programming; 5.3.6 Example 6: Setting Multiple Counters
in a Program\strut
\end{minipage}\tabularnewline
 \begin{minipage}[t]{0.08\columnwidth}\raggedright
16\strut
\end{minipage} & \begin{minipage}[t]{0.30\columnwidth}\raggedright
Can I connect my GL chamber to my main network and access it
remotely?\strut
\end{minipage} & \begin{minipage}[t]{0.33\columnwidth}\raggedright
Your GL chamber can join the main network using DHCP or static. The UI
can be accessed via a Web browser.\strut
\end{minipage} & \begin{minipage}[t]{0.17\columnwidth}\raggedright
8 Network Setup \& Remote Access\strut
\end{minipage}\tabularnewline
 \begin{minipage}[t]{0.08\columnwidth}\raggedright
17\strut
\end{minipage} & \begin{minipage}[t]{0.30\columnwidth}\raggedright
Can I set my GL chamber to use a static IP address on the network?\strut
\end{minipage} & \begin{minipage}[t]{0.33\columnwidth}\raggedright
Your GL chamber can use a static IP address on a static network\strut
\end{minipage} & \begin{minipage}[t]{0.17\columnwidth}\raggedright
8.5 Custom Static Network\strut
\end{minipage}\tabularnewline
 \begin{minipage}[t]{0.08\columnwidth}\raggedright
18\strut
\end{minipage} & \begin{minipage}[t]{0.30\columnwidth}\raggedright
Can I access other GL chambers on the network?\strut
\end{minipage} & \begin{minipage}[t]{0.33\columnwidth}\raggedright
Working with multiple GL chambers on the network, you can choose one GL
controller (on the network) to monitor other GL chambers on the
network.\strut
\end{minipage} & \begin{minipage}[t]{0.17\columnwidth}\raggedright
2.1 GL Controller Operation Options; 9 Network Monitor\strut
\end{minipage}\tabularnewline
 \begin{minipage}[t]{0.08\columnwidth}\raggedright
19\strut
\end{minipage} & \begin{minipage}[t]{0.30\columnwidth}\raggedright
When I issue set command to set a new temp value via TCP communication,
the system responds with PROTECT ON.\strut
\end{minipage} & \begin{minipage}[t]{0.33\columnwidth}\raggedright
If remote TCP communication was set to Monitor only, set command cannot
be issued. Enable remote TCP communication with Monitor and
Control.\strut
\end{minipage} & \begin{minipage}[t]{0.17\columnwidth}\raggedright
2.10.1 Remote Communication Options; 6.5 Set Protection\strut
\end{minipage}\tabularnewline
\begin{minipage}[t]{0.08\columnwidth}\raggedright
20\strut
\end{minipage} & \begin{minipage}[t]{0.30\columnwidth}\raggedright
How do I enter the password to log in on the HMI?\strut
\end{minipage} & \begin{minipage}[t]{0.33\columnwidth}\raggedright
To enter the password to log in on the HMI, the floating input keyboard
must be enabled via the hamburger menu.\strut
\end{minipage} & \begin{minipage}[t]{0.17\columnwidth}\raggedright
2.6.3 How to Log in via HMI\strut
\end{minipage}\tabularnewline
\bottomrule
\end{longtable}